\section{Discussion}
\label{sec:discussion}

\paragraph{Benchmark scores as configuration-conditional measurements.}
The measured pattern---large prompt effects, small backend/quantization
effects, a prompt-induced rank reversal between nearly-equal models, exact
greedy seed-invariance, and high-churn sampling---supports treating a
reported score as a measurement produced by a model \emph{under} an
evaluation configuration, rather than as an intrinsic scalar property of
the model. We state this as an empirical claim about the measured setting
(three small instruction-tuned models, four generative tasks on frozen
200-item slices, one GPU family), not as a universal law; the framework
exists precisely so that others can measure it elsewhere.

\paragraph{What follows for reporting practice.}
Three cheap changes to evaluation reports would have made every claim in
this paper verifiable and most misreadings impossible: (1) publish the
configuration (prompt template, few-shot count, backend and version,
quantization, decoding, seed) next to every number, in machine-readable
form; (2) publish per-item correctness records so that paired tests are
possible, not just aggregate scores; (3) report interval-quantified scores
and, for model comparisons, interval-overlap or paired-test outcomes rather
than point-estimate ranks alone. None of these require new inference: they
require recording what the runner already knows.

\paragraph{Score changes vs.\ ranking changes.}
The concise-prompt result shows that a factor can move scores by 8--17\,\pp{}
without changing any ordering; the five-shot result shows that a factor can
flip the top-two order of models whose control scores differ by only
2.6\,\pp. For model selection, the ranking is the decision-relevant quantity,
and it is precisely the quantity most sensitive to near-ties. Evaluation
reports that compress a model to one number cannot distinguish these
regimes.
\paragraph{Provisional reliability diagnostics.}
As exploratory tooling (not validated methodology), the framework computes
an Evaluation Reliability Score (ERS): a weighted composite of four
components---CI separation among model pairs (0.667 here), seeded bootstrap
ranking stability (mean $\taub=0.84$), within/between config-stability
(0.604), and greedy-seed stability (1.0 when the seed cells are
supplied)---yielding ERS $= 0.711$ (3-component) or $0.740$ (4-component)
on our eight-configuration matrix (appendix \cref{tab:ers}). We present
ERS as a proposed descriptive diagnostic whose components are individually
meaningful, explicitly not as a validated metric; its weights are
provisional and its sensitivity to weight choice can be ablated in the
framework. We also ship (but do not claim results from) infrastructure for
budget-constrained evaluation scheduling, which connects to
efficient-evaluation work such as tinyBenchmarks
\citep{polo2024tinybenchmarks}.

\paragraph{Failure analysis implications.}
The automated categories (\cref{tab:failures}) contain 3 unparseable and
11,175 parseable incorrect responses across 24,800 outputs, with no empty
outputs. This does not imply that parsing is harmless: the selected
extraction mismatch in \cref{tab:failure-analysis} shows that a correct
explicit answer may still be scored incorrectly. Counts of genuine reasoning
errors, demonstration copying, and extraction errors require independent
adjudication and are not inferred from the null-prediction rate.

\section{Threats to Validity and Limitations}
\label{sec:limitations}

We organize threats to validity following conventional categories, adapted
for an artifact-based observational study.

\subsection{Construct validity}
\textit{What we measure may not correspond to what we claim to measure.}

\begin{itemize}
  \item \textbf{Scoring protocol.} Generative exact-match with letter/number
extraction differs from the log-likelihood multiple-choice protocol used by
standard leaderboards.
  \item \textbf{Extractor validation.} The permissive parser has not been
independently adjudicated.
  \item \textbf{Task coverage.} Four tasks do not represent the full
space of LLM capabilities.
  \item \textbf{Contamination.} Benchmark contamination is unmeasured here.
\end{itemize}

\subsection{Internal validity}
\textit{Explanations other than the configuration change may account for observed differences.}

\begin{itemize}
  \item \textbf{Software drift.} Runs span 14 distinct git commits with
varying package versions.
  \item \textbf{Model revisions.} Model and tokenizer hub revisions were
not pinned in run manifests.
  \item \textbf{Precision.} Recorded as \texttt{dtype=auto}.
  \item \textbf{Post-hoc selection.} The Phi/Qwen comparison was
selected after observing top scores.
\end{itemize}

\subsection{External validity}
\textit{Results may not generalize beyond the measured setting.}

\begin{itemize}
  \item \textbf{Model scale.} Three models (1.7B--3.8B) plus one 7B cell.
  \item \textbf{Benchmark slices.} 200-item slices are small.
  \item \textbf{Hardware.} All runs on NVIDIA Tesla T4.
  \item \textbf{Sampling coverage.} Missing Phi-3.5-mini sampling cells.
  \item \textbf{Configuration space.} Other dimensions not explored.
\end{itemize}

\subsection{Statistical conclusion validity}
\textit{Statistical inferences may be weakened by assumption violations.}

\begin{itemize}
  \item \textbf{Item dependence.} MGSM translations and GSM8K questions
are not independent.
  \item \textbf{Multiple testing.} Holm--Bonferroni is conservative.
  \item \textbf{Interval interpretation.} Wilson intervals assume
binomial sampling.
  \item \textbf{Paired test assumptions.} McNemar assumes independent
discordant pairs.
  \item \textbf{Bootstrap scope.} Ranking resamples configurations;
paired differences resample items within tasks. Neither corrects for all
cross-task or translated-item dependence.
\end{itemize}

\subsection{Measurement and recording limitations}
\textit{Data quality and completeness constraints.}

\begin{itemize}
  \item \textbf{Missing cells.} Three Phi-3.5-mini sampling cells missing.
  \item \textbf{Latency placeholders.} vLLM cells store 0.0 timing.
  \item \textbf{Reproducibility scope.} Analysis reproduction feasible.
  \item \textbf{OFAT, not factorial.} Full factorial (Prompt $\times$ Runtime $\times$
Seed) planned and committed to \texttt{t4\_factorial.yaml}.
\end{itemize}

\section{Ethical Considerations}
\label{sec:ethics}
This study uses small, publicly licensed benchmark slices (CC-BY-SA-4.0,
MIT) redistributed at reduced size with attribution and license statements
in the data manifest; users of the slices should consult the originals for
full terms. We did not evaluate safety behaviors, toxicity, or privacy, and
draw no conclusions about them; the repository's 6-item safety canary is a
pipeline check, not a safety evaluation. Benchmark contamination is
unmeasured here (\cref{sec:limitations}), and we caution against reading
absolute accuracies as capability statements. No human subjects data is
involved.

\section{Conclusion}
\label{sec:conclusion}
We presented a configuration-aware, artifact-complete evaluation framework
and used it to measure how three small instruction-tuned LLMs' benchmark
scores and rankings respond to reasonable evaluation-configuration changes
on frozen slices. Prompt configuration was the dominant source of
variability (up to 17.4\,\pp{} within a model, versus a 2.6\,\pp{} top-two
control gap with overlapping intervals); one prompt change produced a
statistically separated reversal of the top-two ranking; quantization and
backend effects were small and not significant after multiple-testing
correction; greedy decoding was exactly seed-invariant; and temperature
sampling induced high per-item churn with small net drift. The practical
conclusion for this setting is to treat a benchmark score as a measurement
conditioned on its configuration: publish the configuration, the per-item
records, and interval-aware comparisons alongside every number. The
framework, registry, and full provenance chain are released so that each
number in this paper can be re-derived, and so that the factorial and
scaling extensions needed to broaden these conclusions can be run on the
same footing.


